@@M0@@

Denoting by @@M1@@ the operations @@M2@@ corresponding to the operations @@M3@@, we therefore have:
\sgaStableTarget{sga6.e05.eq.3-6-2}@@M4@@

We denote by @@M5@@ the coefficient of @@M6@@ in @@M7@@; the reader can verify that it is a polynomial in @@M8@@ whose isobaric components have weights ranging between @@M9@@ and @@M10@@.

Let @@M11@@ be a @@M12@@-ring, @@M13@@. From diagram (3.6.1), we deduce the relations:
\sgaStableTarget{sga6.e05.eq.3-6-3}@@M14@@
\sgaStableTarget{sga6.e05.eq.3-6-4}@@M15@@

Finally, we observe that the pre-@@M16@@-ring structure just defined on @@M17@@ is a @@M18@@-ring structure, since it is obtained from the one defined in 2.3 by the isomorphism @@M19@@.

\sgaStableTarget{sga6.e05.para.3-7}3.7. To avoid confusion, the functor @@M20@@, endowed with the structure of a functor to @@M21@@-rings defined in §2 (resp. in §3), will be denoted by @@M22@@ (resp. @@M23@@). It follows from the preceding considerations that to specify a
